% ===========================================================================
% Section 1: Introduction
%
% Ported from paper/paper_planned.pdf (the authoritative coauthor draft).
% Follows Shapiro's contractual-intro template (see Plan/foursteps.pdf):
%   ¶1-2 Motivation        ¶3-4 Challenges         ¶5 This paper
%   ¶6-7 Model/framework   ¶8-9 Data               ¶10-11 Methods
%   ¶12-13 Findings        ¶14-15 Contributions    ¶16-17 Roadmap
%
% PLACEHOLDERS (all authored by coauthors; preserve as-is):
%   - Paragraph on main findings (pending MA regressions)
%   - Paragraph on counterfactual findings (pending Block 2)
%   - Paragraph on mechanisms findings (pending Block 2)
%   - Numbers in square brackets (pending AutoFill): [10,000], [18,000],
%     [312]
% ===========================================================================

\section{Introduction}
\label{sec:intro}

How do large-scale transport infrastructure changes reshape regional
economies? While most research examines infrastructure expansions
designed to promote development, governments often face a different
challenge: fiscal consolidation that forces difficult choices about
which infrastructure to maintain. Argentina's experience between 1960
and 1991 provides a unique window into this question. Facing severe
fiscal pressure, the government implemented a massive transport
restructuring: closing [10{,}000] kilometers of railway lines while
expanding paved roads by [18{,}000] kilometers. This episode offers
rare evidence on how fiscal consolidation-driven infrastructure
changes—rather than development-oriented investments—affect regional
market access and economic outcomes.

This setting is valuable for two reasons. First, it speaks to a
policy-relevant question: when governments must cut infrastructure
spending, how do these cuts reshape the economic geography? The
Argentine plan, motivated by fiscal rationalization rather than
regional development goals, effectively acted as an unintended
industrial and regional policy by shifting the dominant transport
mode from rail to road. Second, it provides an opportunity to estimate
how regional outcomes respond to market access changes in a multimodal
infrastructure shock. Unlike typical infrastructure studies that
examine single-mode expansions, Argentina's bundled rail-road
restructuring allows us to study elasticities of population and
sectoral employment to market access changes that integrate both
transport modes.

Estimating the causal effects of such infrastructure changes faces
two fundamental challenges. First, infrastructure placement and
removal decisions are endogenous. Even when motivated by fiscal
consolidation, governments may close lines in already-declining
regions or expand roads where growth is expected. This makes it
difficult to separate the causal effect of market access changes
from pre-existing economic trends.

Second, when both rail and road networks change simultaneously, their
effects are difficult to disentangle. The changes are often spatially
correlated—regions losing rail access may or may not gain road
access—creating multicollinearity problems if we try to estimate
mode-specific effects in a single regression. Moreover, the
theoretical framework suggests that what matters for regional
outcomes is total market access integrating all transport modes, not
the separate contribution of each mode. This raises a methodological
question: what can and cannot be identified about mode-specific
effects in multimodal infrastructure episodes?

This paper addresses these challenges by studying Argentina's
transport restructuring through the lens of market access. We
construct a unified market access index that integrates both rail and
road networks using a standard spatial economics framework, computing
generalized transport costs and shortest-path routes on the historical
networks. Our main analysis estimates how regional outcomes—population,
sectoral employment, and economic activity—respond to changes in
total market access induced by the observed infrastructure plan. We
then use counterfactual market access measures to explore the relative
roles of rail closures versus road expansion, constructing separate
scenarios where only rail changes (roads frozen) or only roads change
(rail frozen).

Our empirical framework follows the market access approach standard
in spatial economics and trade. We model each district's economic
outcomes as functions of its market access, defined as the
population-weighted sum of inverse generalized transport costs to all
other districts. To address endogeneity, we instrument for observed
market access changes using two sources of variation. First, we
exploit the Larkin Plan's technical rules: the World Bank study that
motivated the reforms only examined railway segments below a specific
freight threshold for potential closure, creating a discontinuity in
closure probability. Second, we construct hypothetical road networks
connecting major cities via least-cost paths, which predict road
expansion but are orthogonal to local economic shocks.

We construct a novel dataset combining historical transport networks
with census data at the district level, covering [312] districts
across Argentina. We validate our identification strategy with
pre-period balance checks, pre-trend tests, spatial placebo tests,
and over-identification tests.

% [PLACEHOLDER paragraph on counterfactual findings --- fill in after
%  computing counterfactual MA. Expected: rail-driven MA changes
%  account for most effects. Language: "suggestive evidence." Block 2.
%  If we pursue sector-specific analysis, flag that extension here.]

% [PLACEHOLDER paragraph on mechanisms --- fill in after mechanism
%  regressions. Expected: districts losing stations/depots experienced
%  larger declines than MA alone predicts. Block 2.]

Our main results document a contrast between population and sectoral
activity. The elasticity of total population with respect to
$\Delta \ln \mathrm{MA}^{\mathrm{full}}$ is 0.046 (0.033) in the
combined-IV specification; individually, the urban, rural, and
total-population coefficients are not statistically distinguishable
from zero at conventional levels. In contrast, the elasticity of
manufacturing activity is sizable and significant: the log change in
the value of manufacturing production has an elasticity of 0.26
(0.12) and the log change in the manufacturing wage mass has an
elasticity of 0.33 (0.13). Agricultural activity, measured by the
number of farms and by total farmed area, does not respond
detectably. This pattern is consistent with the scale-economies
hypothesis introduced in Section~\ref{sec:history}: the transport
restructuring shifted the cost structure faced by different sectors
differently, with manufacturing benefiting more than agriculture.
Robustness to alternative trade elasticities, alternative
hypothetical-road instruments, and subsample choice is reported in
Section~\ref{sec:results_robustness}.

This paper contributes to several literatures. First, we provide
causal evidence on the regional effects of infrastructure
disinvestment driven by fiscal consolidation. The Argentine episode
involved not just the removal of rail infrastructure but a
technological shift in the dominant transport mode—from rail to
road—that effectively acted as an unintentional industrial policy,
reshaping the country's economic geography. Second, we contribute to
understanding mode-specific effects in multimodal infrastructure
episodes, proposing a counterfactual market access approach that
clarifies what can and cannot be identified. Third, our findings
speak to contemporary policy debates about modal shifts in transport,
as many countries face ongoing transitions between rail and road.
Fourth, we contribute to a longstanding debate in Argentine economic
history about whether railroad closures caused the decline of
hinterland towns or merely followed pre-existing trends.

The paper proceeds as follows. Section~\ref{sec:history} provides
historical context. Section~\ref{sec:data} describes data and market
access construction. Section~\ref{sec:identification} presents the
empirical strategy. Section~\ref{sec:results} reports main results.
Section~\ref{sec:counterfactuals} explores mode-specific patterns.
Section~\ref{sec:mechanisms} examines mechanisms and heterogeneity.
Section~\ref{sec:conclusion} concludes.
